\documentclass[12pt,a4paper]{article}
\usepackage[utf8]{inputenc}
\usepackage[french]{babel}
\usepackage[T1]{fontenc}
\usepackage{amsmath}
\usepackage{amsfonts}
\usepackage{amssymb}
\usepackage{graphicx}
\usepackage{hyperref}
\usepackage{titlesec}
\usepackage{geometry}
\usepackage{listings}
\usepackage{color}

\geometry{margin=2.5cm}

\definecolor{codegray}{rgb}{0.5,0.5,0.5}
\definecolor{codegreen}{rgb}{0,0.6,0}
\definecolor{codeblue}{rgb}{0,0,0.8}

\lstset{
    backgroundcolor=\color{white},
    commentstyle=\color{codegreen},
    keywordstyle=\color{codeblue},
    numberstyle=\tiny\color{codegray},
    stringstyle=\color{red},
    basicstyle=\ttfamily\footnotesize,
    breaklines=true,
    numbers=left,
    frame=single
}

\title{Recapitulatif Technique : Structuration G\'en\'erique par Pays\\ \large Projet Mon \'Ecole}
\author{D\'epartement Technique - MonEkole}
\date{21 Janvier 2026}

\begin{document}

\maketitle

\section{Objectif du Projet}
L'objectif principal est de rendre le système \textbf{Mon \'Ecole} capable de s'adapter à n'importe quel système éducatif national sans modification du code source (modélisation dynamique). La flexibilité se concentre sur la hiérarchie organisationnelle entre le niveau provincial et l'établissement scolaire.

\section{Logique de la Nouvelle Base de Données}

\subsection{Concept de ``Steps'' (Niveaux)}
Chaque pays possède une profondeur hiérarchique différente. La logique de la base \texttt{countryStructure} repose sur deux axes :
\begin{itemize}
    \item \textbf{Structuration P\'edagogique} : Organisation des entités d'enseignement (ex: DPE $\to$ DCE $\to$ Zone $\to$ \'Ecole).
    \item \textbf{Structuration Administrative} : Organisation territoriale et géographique (ex: Province $\to$ Commune $\to$ Secteur $\to$ Colline).
\end{itemize}

\subsection{Isolation et S\'ecurit\'e}
Contrairement aux données opérationnelles (élèves, notes, finances), la configuration de la structure nationale est stockée dans une base de données isolée nommée \textbf{\texttt{countryStructure}}. Cette séparation permet :
\begin{enumerate}
    \item Une maintenance facilitée sans risque de corruption des données scolaires.
    \item Une portabilité de la configuration d'un pays à l'autre.
    \item Un routage spécifique au niveau applicatif via Django.
\end{enumerate}

\section{\'Etapes de R\'ealisation}

\subsection{Infrastructure de Donn\'ees}
\begin{itemize}
    \item \textbf{Cr\'eation de la Base} : Mise en place de \texttt{countryStructure} sur le serveur MySQL.
    \item \textbf{Mise en place de la Nomenclature} :
    \begin{itemize}
        \item Table \texttt{pays} : Stockage du nombre de niveaux ($n$ Pedagogique, $m$ Administratif).
        \item Table \texttt{structuresPedagogiques} : Liste ordonnée des noms de niveaux et codes.
        \item Table \texttt{structuresAdministratives} : Idem pour l'axe territorial.
    \end{itemize}
\end{itemize}

\subsection{D\'eveloppement Backend (Django)}
\begin{itemize}
    \item \textbf{Multi-Database Routing} : Création d'un \texttt{Database Router} pour intercepter les requêtes vers les modèles de structuration et les diriger vers la base secondaire.
    \item \textbf{APIs REST} : Développement d'un endpoint pour la sauvegarde, la récupération et l'ajustement dynamique des niveaux scolaires.
    \item \textbf{Logique de Code} : G\'en\'eration automatique de codes sur 3 caractères (ex: Province $\to$ PRO) avec gestion des accents.
\end{itemize}

\subsection{Interface Utilisateur (UI/UX)}
\begin{itemize}
    \item \textbf{Dashboard Moderne} : Utilisation du Glassmorphism, de gradients professionnels et de micro-interactions.
    \item \textbf{R\`egles de Coh\'erence} :
    \begin{itemize}
        \item Validation en temps réel du nombre d'éléments saisis par rapport aux limites définies.
        \item \textbf{Protection anti-quitter} : JavaScript empêche la fermeture accidentelle de la page si la structure n'est pas complète.
        \item \textbf{Ajustement Intelligent} : Option pour recalibrer automatiquement les compteurs de niveaux en fonction des saisies réelles.
    \end{itemize}
\end{itemize}

\section{Architecture Technique (Extrait Code)}
Le routage est géré de la manière suivante dans Django :
\begin{lstlisting}[language=Python]
class CountryStructureRouter:
    def db_for_read(self, model, **hints):
        if model.__name__ in ['Pays', 'StructurePedagogique', 'StructureAdministrative']:
            return 'countryStructure'
        return 'default'
\end{lstlisting}

\section{Conclusion}
Le système est d\`esormais pr\^et \`a accueillir n'importe quel pays de la zone Afrique de l'Est (et au-del\`a). La structure est totalement d\'ecoupl\'ee du reste de l'application, garantissant robustesse et scalabilit\'e.

\end{document}
